
Economic recessions represent periods of significant contraction in economic activity that impose substantial costs on society through job losses, reduced output, and diminished welfare. Accurately forecasting these downturns remains a central challenge in macroeconomics and financial economics, with important implications for policymakers, investors, and the broader public \citep{murphyPredictingRecessions1959}. Among the various indicators examined in the literature, the yield curve, the term structure of interest rates across different maturities, has emerged as one of the most reliable predictors of future economic contractions.

The predictive power of the yield curve, particularly the spread between long-term and short-term Treasury yields, has been well documented since the seminal work of \citet{estrellaYieldCurvePredictor1996}. Their probit-based framework demonstrated that inversions of the yield curve, where short-term rates exceed long-term rates, reliably signal recessions four to six quarters ahead. The economic intuition underlying this relationship is compelling: an inverted curve reflects market expectations of future monetary policy easing in response to anticipated economic weakness, as investors price in lower future short-term rates relative to current levels. Recent work by \citet{yorkPredictingRecessionProbabilities2018} and \citet{sharpeNearTermForwardYield2018} has refined this approach by focusing on near-term forward spreads, which more directly capture market expectations of monetary policy trajectories over the next 12-18 months.

However, the post-2008 financial landscape has been characterized by unprecedented monetary policy interventions, including quantitative easing programs, forward guidance, and extended periods at the zero lower bound. These unconventional policies have fundamentally altered the behavior of long-term yields by compressing term premia, the additional compensation investors demand for holding longer-maturity bonds. As \citet{rudebusch2006macroeconomic} emphasize, movements in long-term rates embody both changes in expectations about future short rates and variations in the term premium, and these two components can have markedly different macroeconomic implications. During periods of aggressive asset purchases by central banks, term premium compression can flatten or even invert the yield curve without necessarily signaling expectations of future rate cuts, thereby potentially generating false recession signals. \citet{bauer2018information} argue that adjusting for term premia does not necessarily improve forecast accuracy in model-free settings, though they acknowledge the importance of distinguishing these components for structural interpretation.

Moreover, the COVID-19 pandemic introduced an entirely new type of economic shock, an exogenous public health crisis that triggered sudden stops in economic activity largely independent of financial market conditions. This episode raised fundamental questions about whether traditional financial indicators, including yield curve factors, retain their predictive power during non-financial crises. The sharp divergence between financial signals and real economic outcomes during 2020 suggests that recession forecasting models must be adaptive to different crisis regimes.

This paper makes several contributions to the literature on yield curve-based recession forecasting. First, we extend the traditional spread-based approach by decomposing the entire yield curve into three latent factors, Level, Slope, and Curvature, using Principal Component Analysis rather than relying solely on parametric specifications such as the \citet{nelsonParsimoniousModelingYield1987} framework or the dynamic extension proposed by \citet{diebold2006forecasting}. This data-driven approach captures the full dimensionality of term structure movements and allows us to examine how each factor's relationship with future recessions varies across regimes. Second, we explicitly address the term premium problem identified by \citet{rudebusch2006macroeconomic} by incorporating Kim-Wright decompositions to construct adjusted slope measures that isolate expectations-driven yield movements from non-expectational forces. Third, we move beyond linear probit models to employ modern machine learning techniques, including Ridge Regression, Elastic Net, Random Forests, and Gradient Boosting, that can flexibly capture nonlinear relationships and structural breaks. Fourth, we implement a rigorous out-of-sample evaluation framework with strict time-respecting procedures to avoid look-ahead bias and assess true real-time forecasting performance.

Our empirical analysis spans from 1982 to 2025 and includes detailed scenario studies of the two most recent major recessions: the Global Financial Crisis and the COVID-19 recession. We construct multiple feature sets combining yield curve factors with real-sector macroeconomic indicators such as unemployment, industrial production, and inflation to test whether hybrid models can adapt across different shock types. The results reveal striking regime dependence: during the GFC, yield PCA factors dominate with balanced accuracy exceeding 0.83, while during COVID-19, macro-only models achieve 0.86 balanced accuracy as yield signals fail almost entirely. Feature importance analysis confirms that the slope factor (PC2) and traditional spreads are most informative during financial crises, whereas the level (PC1) and curvature (PC3) factors gain prominence during periods of intense policy intervention.

These findings have important practical implications. For policymakers and market participants, our results suggest that single-indicator models are insufficient in the post-crisis era. Instead, robust recession forecasting requires ensemble approaches that dynamically weight financial versus real-sector signals based on the prevailing economic regime. We conclude by discussing potential extensions, including regime-switching meta-models, integration of high-frequency data such as weekly jobless claims, and further decomposition of yield factors into risk-neutral and term premium components.
